% Options for packages loaded elsewhere
\PassOptionsToPackage{unicode}{hyperref}
\PassOptionsToPackage{hyphens}{url}
\PassOptionsToPackage{space}{xeCJK}
%
\documentclass[
  ignorenonframetext,
  aspectratio=169,
]{beamer}
\title{相互作用行列のネットワーク解析}
\subtitle{推定した相互作用行列をグラフとして読む ---
keystone・bridge・栄養層・リワイヤリング}
\author{西岡佳祐 --- NIFE / SFB TRR-298}
\date{2026-06-03}

\usepackage{pgfpages}
\setbeamertemplate{caption}[numbered]
\setbeamertemplate{caption label separator}{: }
\setbeamercolor{caption name}{fg=normal text.fg}
\beamertemplatenavigationsymbolsempty
% Prevent slide breaks in the middle of a paragraph
\widowpenalties 1 10000
\raggedbottom
\setbeamertemplate{part page}{
  \centering
  \begin{beamercolorbox}[sep=16pt,center]{part title}
    \usebeamerfont{part title}\insertpart\par
  \end{beamercolorbox}
}
\setbeamertemplate{section page}{
  \centering
  \begin{beamercolorbox}[sep=12pt,center]{part title}
    \usebeamerfont{section title}\insertsection\par
  \end{beamercolorbox}
}
\setbeamertemplate{subsection page}{
  \centering
  \begin{beamercolorbox}[sep=8pt,center]{part title}
    \usebeamerfont{subsection title}\insertsubsection\par
  \end{beamercolorbox}
}
\AtBeginPart{
  \frame{\partpage}
}
\AtBeginSection{
  \ifbibliography
  \else
    \frame{\sectionpage}
  \fi
}
\AtBeginSubsection{
  \frame{\subsectionpage}
}
\usepackage{amsmath,amssymb}
\usepackage{lmodern}
\usepackage{iftex}
\ifPDFTeX
  \usepackage[T1]{fontenc}
  \usepackage[utf8]{inputenc}
  \usepackage{textcomp} % provide euro and other symbols
\else % if luatex or xetex
  \usepackage{unicode-math}
  \defaultfontfeatures{Scale=MatchLowercase}
  \defaultfontfeatures[\rmfamily]{Ligatures=TeX,Scale=1}
  \setmainfont[]{Noto Serif CJK JP}
  \setmonofont[]{Noto Sans Mono CJK JP}
  \ifXeTeX
    \usepackage{xeCJK}
    \setCJKmainfont[]{Noto Serif CJK JP}
  \fi
  \ifLuaTeX
    \usepackage[]{luatexja-fontspec}
    \setmainjfont[]{Noto Serif CJK JP}
  \fi
\fi
\usetheme[]{Madrid}
\usecolortheme{whale}
\usefonttheme{serif} % use mainfont rather than sansfont for slide text
% Use upquote if available, for straight quotes in verbatim environments
\IfFileExists{upquote.sty}{\usepackage{upquote}}{}
\IfFileExists{microtype.sty}{% use microtype if available
  \usepackage[]{microtype}
  \UseMicrotypeSet[protrusion]{basicmath} % disable protrusion for tt fonts
}{}
\makeatletter
\@ifundefined{KOMAClassName}{% if non-KOMA class
  \IfFileExists{parskip.sty}{%
    \usepackage{parskip}
  }{% else
    \setlength{\parindent}{0pt}
    \setlength{\parskip}{6pt plus 2pt minus 1pt}}
}{% if KOMA class
  \KOMAoptions{parskip=half}}
\makeatother
\usepackage{xcolor}
\IfFileExists{xurl.sty}{\usepackage{xurl}}{} % add URL line breaks if available
\IfFileExists{bookmark.sty}{\usepackage{bookmark}}{\usepackage{hyperref}}
\hypersetup{
  pdftitle={相互作用行列のネットワーク解析},
  pdfauthor={西岡佳祐 --- NIFE / SFB TRR-298},
  hidelinks,
  pdfcreator={LaTeX via pandoc}}
\urlstyle{same} % disable monospaced font for URLs
\newif\ifbibliography
\usepackage{graphicx}
\makeatletter
\def\maxwidth{\ifdim\Gin@nat@width>\linewidth\linewidth\else\Gin@nat@width\fi}
\def\maxheight{\ifdim\Gin@nat@height>\textheight\textheight\else\Gin@nat@height\fi}
\makeatother
% Scale images if necessary, so that they will not overflow the page
% margins by default, and it is still possible to overwrite the defaults
% using explicit options in \includegraphics[width, height, ...]{}
\setkeys{Gin}{width=\maxwidth,height=\maxheight,keepaspectratio}
% Set default figure placement to htbp
\makeatletter
\def\fps@figure{htbp}
\makeatother
\setlength{\emergencystretch}{3em} % prevent overfull lines
\providecommand{\tightlist}{%
  \setlength{\itemsep}{0pt}\setlength{\parskip}{0pt}}
\setcounter{secnumdepth}{-\maxdimen} % remove section numbering
\usepackage{amsmath,amssymb}
\newcommand{\sgn}{\operatorname{sgn}}
\newcommand{\relu}[1]{\left[#1\right]_{+}}
\ifLuaTeX
  \usepackage{selnolig}  % disable illegal ligatures
\fi

\begin{document}
\frame{\titlepage}

\begin{frame}{このデッキの位置づけ}
\protect\hypertarget{ux3053ux306eux30c7ux30c3ux30adux306eux4f4dux7f6eux3065ux3051}{}
本プロジェクト = 3 本柱 ＋ 空間拡張。データの流れ：

raw 16S → guild \(\varphi\) → gLV/Hamilton（＋符号 prior）→ LOO 検証 →
空間 PDE

デッキ構成（◀ = 本デッキ）：

\begin{itemize}
\tightlist
\item
  \textbf{Overview}（傘）--- 全体像と3本柱
\item
  \textbf{AGORA} --- 代謝 → 符号 prior（生態モデルの入力）
\item
  \textbf{Dieckow} --- in-vivo 相互作用の推定と検証（生態モデル本体）
\item
  \textbf{Network} --- 相互作用行列 \(A\) の構造解析 ◀ \textbf{本デッキ}
\item
  \textbf{Spatial-PDE} --- FISH 深さプロファイルの反応拡散
\item
  \textbf{FISH pipeline} --- .lif → 5 種深さ組成
\end{itemize}
\end{frame}

\begin{frame}{動機：\(A\) を「当てはめる」から「読む」へ}
\protect\hypertarget{ux52d5ux6a5fa-ux3092ux5f53ux3066ux306fux3081ux308bux304bux3089ux8aadux3080ux3078}{}
相互作用行列 \(A\) を推定するだけでなく、その\textbf{グラフ構造}を読む。

\begin{itemize}
\tightlist
\item
  \textbf{keystone}：少量でも群集を支配する種（高い中心性）。
\item
  \textbf{bridge}：栄養層をつなぐ媒介種（高い betweenness）。
\item
  \textbf{trophic layer}：cross-feeding の生産者 \(\to\)
  消費者の階層性。
\item
  \textbf{rewiring}：健常 \(\leftrightarrow\) dysbiosis
  でエッジ符号が反転する再配線。
\end{itemize}

\vspace{0.4em}

口腔バイオフィルムの古典像は「\emph{P. gingivalis} = keystone、\emph{F.
nucleatum} = bridge」。 本デッキは、in-vivo class
レベルでこの像が\textbf{支持されるか}をネットワーク統計で検証する。
\end{frame}

\begin{frame}{ゲージ不変性：\(A\) ではなく \(A_{\text{eff}}\) を読む}
\protect\hypertarget{ux30b2ux30fcux30b8ux4e0dux5909ux6027a-ux3067ux306fux306aux304f-a_texteff-ux3092ux8aadux3080}{}
replicator 力学
\[\dot{\varphi}_i = \varphi_i\Big[(A\varphi+b)_i - \varphi^{\!\top}(A\varphi+b)\Big],\qquad \sum_i\varphi_i=1\]
は \(A\) の\textbf{列ごとの定数シフト} \(A_{ij}\to A_{ij}+c_j\)
に対して不変（\(\sum_i\varphi_i=1\) で消える）。

\(\Rightarrow\) 生の \(A_{ij}\)
の符号・大きさは\textbf{ゲージ依存}で意味を持たない。\textbf{列中心化}した
\[A_{\text{eff}}[i,j] = A_{ij} - \frac{1}{S}\sum_{k} A_{kj}\]
のみがゲージ不変。\textbf{以降の全ての符号・中心性解釈は
\(A_{\text{eff}}\) に対して行う。}

有向重み付きグラフ \(G=(V,E)\) を \(w_{ij}=A_{\text{eff}}[i,j]\)
で定義する（\(|V|=S=10\) ギルド）。
\end{frame}

\begin{frame}{中心性の定義}
\protect\hypertarget{ux4e2dux5fc3ux6027ux306eux5b9aux7fa9}{}
\(A_{\text{eff}}\) から複数の中心性を計算：

\begin{itemize}
\tightlist
\item
  \textbf{固有ベクトル中心性}：\(A\,v = \lambda\,v\) の主固有ベクトル
  \(v\)。 \(v_i\)
  は「中心的な相手とつながる種ほど中心的」を再帰的に表す。
\item
  \textbf{媒介中心性 (betweenness)}：最短経路が頂点 \(i\) を通る割合
  \[g(i)=\sum_{s\neq i\neq t}\frac{\sigma_{st}(i)}{\sigma_{st}}.\]
\item
  \textbf{PageRank}：定常分布
  \(\pi = (1-d)\,\mathbf{1}/S + d\,M\pi\)（\(M\) は列正規化推移行列）。
\item
  \textbf{入次/出次強度}：\(\;s^{\text{in}}_i=\sum_j |A_{\text{eff}}[i,j]|,\quad s^{\text{out}}_j=\sum_i |A_{\text{eff}}[i,j]|.\)
\end{itemize}
\end{frame}

\begin{frame}{中心性の結果（Dieckow 10 ギルド・no-prior consensus）}
\protect\hypertarget{ux4e2dux5fc3ux6027ux306eux7d50ux679cdieckow-10-ux30aeux30ebux30c9no-prior-consensus}{}
\includegraphics[width=\textwidth,height=0.5\textheight]{results/guild_network/guild_centrality_summary.png}

固有ベクトル中心性：Bacilli \textbf{0.61}、Actinobacteria
\textbf{0.58}、Betaproteobacteria 0.44、 Negativicutes
0.22、Bacteroidia(\(\approx\) \emph{P. gingivalis}) 0.19。 媒介：Bacilli
\textbf{0.78}、Actinobacteria 0.22、他は 0。 平均存在量：Bacilli
\textbf{0.547}（優占）、Actinobacteria 0.203、Bacteroidia 0.047。
\end{frame}

\begin{frame}{keystone / bridge 検定：古典像は成り立つか}
\protect\hypertarget{keystone-bridge-ux691cux5b9aux53e4ux5178ux50cfux306fux6210ux308aux7acbux3064ux304b}{}
中心性と存在量のランクを照合：

\begin{itemize}
\tightlist
\item
  \emph{P. gingivalis}（Bacteroidia）は固有ベクトル中心性ランク
  \(\sim\)\textbf{4}、存在量ランク \(\sim\)\textbf{5}。 \(\Rightarrow\)
  \textbf{構造的 keystone ではない}（少量かつ低中心性）。
\item
  bridge（最高 betweenness）は
  \textbf{Bacilli（\emph{Streptococcus}）}であり、\textbf{Fusobacterium
  ではない}。
\end{itemize}

\vspace{0.3em}

\(\Rightarrow\) 古典的「Pg-keystone / Fn-bridge」像は \textbf{in-vivo
class レベルでは支持されない}。

\vspace{0.3em}

\textcolor{red}{注意（過大解釈を避ける）：class $\neq$ 種；密なグラフは betweenness を縮退させる；
希少ギルドの推定はノイズが大きい。}
\end{frame}

\begin{frame}{影響力 vs 脆弱性（出次強度 vs 入次強度）}
\protect\hypertarget{ux5f71ux97ffux529b-vs-ux8106ux5f31ux6027ux51faux6b21ux5f37ux5ea6-vs-ux5165ux6b21ux5f37ux5ea6}{}
\includegraphics[width=\textwidth,height=0.55\textheight]{results/guild_network/influence_vulnerability.png}

出次強度 \(s^{\text{out}}_j=\sum_i|A_{\text{eff}}[i,j]|\) は \(j\)
が他者に及ぼす\textbf{影響力}、 入次強度
\(s^{\text{in}}_i=\sum_j|A_{\text{eff}}[i,j]|\) は \(i\)
の\textbf{脆弱性}を測る。 両者の非対称性が「駆動側／被駆動側」を分ける。
\end{frame}

\begin{frame}{栄養的コヒーレンス（有向 AGORA cross-feeding）}
\protect\hypertarget{ux6804ux990aux7684ux30b3ux30d2ux30fcux30ecux30f3ux30b9ux6709ux5411-agora-cross-feeding}{}
有向 cross-feeding グラフに対し、MacKay の\textbf{栄養的非コヒーレンス}
\(F_0\) を測る （栄養レベル \(h_i\)
を割り当て、エッジ両端のレベル差の分散）：
\[F_0 = \frac{1}{w}\sum_{ij} W_{ij}\,(h_i - h_j - 1)^2 .\]

\begin{itemize}
\tightlist
\item
  実測 \(F_0 = \mathbf{0.652}\) vs ランダム \(0.646\)、\(p=0.50\)。
\item
  \(\Rightarrow\) \textbf{栄養的にインコヒーレント} ＝
  階層を持たず、\textbf{相互 cross-feeding の循環}。
\end{itemize}

\vspace{0.3em}

生物学的に妥当：Fusobacterium は基底生産者、Veillonella
は高位消費者だが、
ループ状に絡み合い、明確な一方向の食物連鎖を成さない。
\end{frame}

\begin{frame}{Concordant backbone：生態層 \(\times\) 代謝層の合意}
\protect\hypertarget{concordant-backboneux751fux614bux5c64-times-ux4ee3ux8b1dux5c64ux306eux5408ux610f}{}
独立な 2 層を重ねる --- 生態(Hamilton)層と AGORA 代謝層 ---
で両者が\textbf{符号一致}するエッジを抽出：

\begin{itemize}
\tightlist
\item
  \textbf{11 エッジ、すべて正（協調）}、\(p=4\times10^{-4}\)（AGORA
  デッキと同一の独立検証）。
\item
  Negativicutes（\emph{Veillonella}）が主要な\textbf{代謝シンク}：
  \[\{\text{Bacilli, Bacteroidia, }\beta\text{Proteo, Fusobacteria}\}\ \longrightarrow\ \text{Negativicutes}.\]
\item
  これは再構成された\textbf{乳酸 cross-feeding} に対応。
\end{itemize}

\vspace{0.3em}

\textbf{本ネットワーク解析で最も頑健な結果。}符号レベルの backbone
は両層から独立に支持される。
\end{frame}

\begin{frame}{ネットワークの LOO 安定性}
\protect\hypertarget{ux30cdux30c3ux30c8ux30efux30fcux30afux306e-loo-ux5b89ux5b9aux6027}{}
\includegraphics[width=\textwidth,height=0.55\textheight]{results/guild_network/loo_stability.png}

患者を 1 人ずつ抜く leave-one-out で \(A_{\text{eff}}\)
を再推定し、各ペアの符号一致率を測る。 \textbf{全ペアで符号一致率
\(\geq 0.70\)} \(\Rightarrow\)
ネットワークの符号構造は個々の患者に頑健。
\end{frame}

\begin{frame}{エッジ強度の permutation 検定（誠実に）}
\protect\hypertarget{ux30a8ux30c3ux30b8ux5f37ux5ea6ux306e-permutation-ux691cux5b9aux8aa0ux5b9fux306b}{}
\includegraphics[width=\textwidth,height=0.48\textheight]{results/guild_network/permutation_test.png}

各 off-diagonal エッジの\textbf{大きさ}を permutation null と比較：
\textbf{90 ペア中わずか 2 ペア}が個別に \(p<0.05\) を超えるのみ。

\vspace{0.3em}

\(\Rightarrow\)
ネットワークは密で、個々のエッジ\textbf{強度は弱くしか分解されない}。
頑健なのは前スライドの\textbf{符号レベル backbone}であって、個別
magnitude ではない。
\end{frame}

\begin{frame}{CS \(\leftrightarrow\) DH リワイヤリング（Heine 5
種・posterior）}
\protect\hypertarget{cs-leftrightarrow-dh-ux30eaux30efux30a4ux30e4ux30eaux30f3ux30b0heine-5-ux7a2eposterior}{}
5 種アトラクター（CS = commensal-static、DH = dysbiotic-HOBIC）の
posterior を、 列中心化 \(A_{\text{eff}}\)
で比較する。\textbf{促進確率}を
\[P_{\text{facilitation}}(i\!\leftarrow\! j) = \Pr\big[\,A_{\text{eff}}[i,j] > 0\,\big]\]
（posterior 上で \(A_{\text{eff}}[i,j]>0\) となる確率）と定義。

\(P_{\text{fac}}\to 1\) は相互作用が\textbf{促進}、\(\to 0\)
は\textbf{競争／抑制}を意味する。 CS と DH の間で \(P_{\text{fac}}\)
がどう動くかが「再配線」を定量化する。
\end{frame}

\begin{frame}{リワイヤリングの結果}
\protect\hypertarget{ux30eaux30efux30a4ux30e4ux30eaux30f3ux30b0ux306eux7d50ux679c}{}
CS \(\to\) DH で促進確率 \(P_{\text{fac}}\) が劇的に変化：

\begin{itemize}
\tightlist
\item
  \textbf{So \(\leftrightarrow\) Vd
  相互扶助の崩壊}：\emph{Streptococcus}--\emph{Veillonella} 共生が
  dysbiosis で競争へ。
  \[P_{\text{fac}}(\text{So}\!\to\!\text{Vd}):\ 1.00 \to 0.03,\qquad
    P_{\text{fac}}(\text{Vd}\!\to\!\text{So}):\ 0.9995 \to 0.00.\]
\item
  \textbf{Fusobacterium が支えられる}：An\(\to\)Fn と Vd\(\to\)Fn が競争
  \(\to\) 促進へ反転。
  \[P_{\text{fac}}(\text{An}\!\to\!\text{Fn}):\ 0.23 \to 0.97,\qquad
    P_{\text{fac}}(\text{Vd}\!\to\!\text{Fn}):\ 0.09 \to 0.77.\]
\end{itemize}
\end{frame}

\begin{frame}{\emph{P. gingivalis} の中心化}
\protect\hypertarget{p.-gingivalis-ux306eux4e2dux5fc3ux5316}{}
固有ベクトル中心性：CS \textbf{0.32} \(\to\) DH \textbf{0.51}（DH
で最高）。

\begin{itemize}
\tightlist
\item
  ただし DH で Pg は\textbf{正味抑制的}：出力エッジの
  \(P_{\text{facilitation}} = 0.27\)。
\item
  \(\Rightarrow\) Pg は\textbf{静的 keystone ではない}が、動的な
  dysbiotic 状態で \textbf{keystone 様}に「なる」。
\end{itemize}

\vspace{0.3em}

これは FISH の知見と整合：dysbiosis で Pg は深部に沈み、近傍依存を脱して
\textbf{自律化（中心化）}する。keystone
は状態の関数であり、固定属性ではない。
\end{frame}

\begin{frame}{結論}
\protect\hypertarget{ux7d50ux8ad6}{}
\begin{enumerate}
\tightlist
\item
  \textbf{代謝 backbone は頑健}：Veillonella を主シンクとする乳酸
  cross-feeding が、 生態層・代謝層から独立に符号支持される（11
  エッジ・全正・\(p=4\times10^{-4}\)）。
\item
  \textbf{教科書的 keystone は支持されない}：in-vivo class レベルで Pg
  は構造的 keystone でなく、 bridge は Fusobacterium でなく Bacilli。
\item
  \textbf{dysbiosis は abundance シフトでなく「リワイヤリング」}：
  So--Vd 相互扶助の喪失 ＋ Pg の中心化。
\item
  個々のエッジ強度は弱く分解される（permutation で 2/90）---
  頑健なのは符号構造。 class \(\neq\)
  種・密グラフの縮退・希少ギルドのノイズは引き続き限界。
\end{enumerate}
\end{frame}

\end{document}
